\documentclass[12pt]{article}

% ---------- Core packages ----------
\usepackage[margin=1in]{geometry}
\usepackage[T1]{fontenc}
\usepackage{lmodern}
\usepackage{amsmath,amssymb,amsthm}
\usepackage{graphicx}
\usepackage{booktabs}
\usepackage{natbib}
\usepackage{hyperref}
\usepackage{setspace}
\usepackage{lineno}

\hypersetup{colorlinks=true, linkcolor=blue, citecolor=blue, urlcolor=blue}
\graphicspath{{./figures/}}

% ---------- Review-draft line numbers ----------
\linenumbers

% Patches so lineno works with amsmath environments.
\newcommand*\patchAmsMathEnvironmentForLineno[1]{%
  \expandafter\let\csname old#1\expandafter\endcsname\csname #1\endcsname
  \expandafter\let\csname oldend#1\expandafter\endcsname\csname end#1\endcsname
  \renewenvironment{#1}{\linenomath\csname old#1\endcsname}{\csname oldend#1\endcsname\endlinenomath}%
}
\newcommand*\patchBothAmsMathEnvironmentsForLineno[1]{%
  \patchAmsMathEnvironmentForLineno{#1}%
  \patchAmsMathEnvironmentForLineno{#1*}%
}
\AtBeginDocument{%
  \patchBothAmsMathEnvironmentsForLineno{equation}%
  \patchBothAmsMathEnvironmentsForLineno{align}%
  \patchBothAmsMathEnvironmentsForLineno{gather}%
  \patchBothAmsMathEnvironmentsForLineno{multline}%
}

\title{<Manuscript Title>}
\author{<Author Name(s)>}
\date{\today}

\begin{document}
\maketitle

\begin{abstract}
% Default target 6-8 sentences (acceptable 4-10).
% No citations and no math notation in abstract text.
<Write abstract here.>
\end{abstract}

\noindent\textbf{Keywords:} <keyword 1>; <keyword 2>; <keyword 3>; <keyword 4>

\doublespacing

\section{Introduction}\label{sec:introduction}

% Answer three questions:
% (1) Why is this important?
% (2) What has been done?
% (3) What is new here?

Textual citation example: \citet{smith2020example}.
Parenthetical citation example: \citep{smith2020example}.

Roadmap example: Section~\ref{sec:data} describes data,
Section~\ref{sec:methods} presents methods,
Section~\ref{sec:results} reports results, and
Section~\ref{sec:discussion} concludes.

\section{Data}\label{sec:data}

Describe source, population, sample size, variables, and missingness handling.

\section{Methods}\label{sec:methods}

State estimand, assumptions, model/procedure, and uncertainty quantification.

\begin{equation}\label{eq:model}
Y_i = X_i^\top \beta + \varepsilon_i.
\end{equation}

Equation~\eqref{eq:model} defines the baseline model.

\section{Results}\label{sec:results}

Reference each figure/table and interpret findings in text.

\begin{table}[tbp]
\centering
\caption{Example table caption.}
\label{tab:example}
\begin{tabular}{lr}
\toprule
Method & Value \\
\midrule
A & 1.23 \\
B & 1.45 \\
\bottomrule
\end{tabular}
\end{table}

\begin{figure}[tbp]
\centering
% Prefer vector graphics (.pdf/.eps) for line graphics.
\fbox{\rule{0pt}{1.5in}\rule{0.8\linewidth}{0pt}}
\caption{Example figure caption.}
\label{fig:example}
\end{figure}

Figure~\ref{fig:example} and Table~\ref{tab:example} summarize key results.

\section{Discussion}\label{sec:discussion}

Summarize contributions, limitations, and future work.

\section*{Reproducibility and Supplement}

\textbf{Code availability:} <repository or statement>\\
\textbf{Data availability:} <location or access conditions>\\
\textbf{Supplement:} <appendix/supplement file description>

% Uncomment and update for bibliography.
% \bibliographystyle{plainnat}
% \bibliography{refs}

\end{document}
